\section{Integration with the Four-Loop Self-Improvement Architecture}

\label{sec:integration}

The Hanzo Dev agent implements a four-loop self-improvement architecture
inspired by~\cite{shinn2023reflexion} and extended with proactive tool
creation~\cite{wang2023voyager}. Agent-GRPO integrates with each loop.

\subsection{Loop 0: Telemetry}

Loop 0 instruments every tool call with structured telemetry:
\begin{equation}
    \text{TelemetryEvent} = \bigl(\texttt{session\_id},\; t_i,\; \bm{p}_i,\; o_i,\; s_i,\; \Delta t_i,\; \text{context}\bigr),
    \label{eq:telemetry}
\end{equation}
where $\Delta t_i$ is the wall-clock duration of the tool call and \text{context}
captures the surrounding conversational context.

Telemetry events are aggregated at the end of each session to construct the
raw trajectory $\tau$ according to Definition~\ref{def:trajectory}. Loop 0
thus provides the \emph{raw data} for Agent-GRPO's trajectory processing
pipeline without any additional instrumentation cost.

\subsection{Loop 1: Build-It-Now}

Loop 1 monitors for repeated tool-call failures and proactively proposes
new tool implementations~\cite{wang2023voyager}. Agent-GRPO informs this loop
via two mechanisms:

\begin{enumerate}[leftmargin=1.5em]
    \item \textbf{Failure pattern retrieval:} When Loop 1 detects a recurring
          failure (e.g., the same error message appearing across three or more
          sessions), it queries $\mathcal{E}$ for failure pattern experiences
          matching the error. If no remediation experience exists, Loop 1
          escalates to tool creation.
    \item \textbf{Tool selection experience audit:} Experience entries of the
          form ``prefer \texttt{X} over \texttt{Y}'' indicate that tool \texttt{X}
          satisfies a need better than \texttt{Y}. If \texttt{X} does not exist
          in the tool registry $\mathcal{T}$, Loop 1 receives a tool creation
          proposal.
\end{enumerate}

\subsection{Loop 2: Active Learning}

Loop 2 captures explicit user corrections as high-confidence, high-priority
experiences. When a user corrects the agent's output (triggering
$\rho_a(\tau) = 0$), the correction is processed by a dedicated prompt:

\begin{quote}
\textit{The user corrected the agent's response. Original task: \{x\}. Agent
output: \{y\}. User correction: \{y'\}. In $\leq$32 words, what should the
agent have done differently? Format as an actionable guideline.}
\end{quote}

The resulting experience is added to $\mathcal{E}$ with $c = 0.8$ (high
confidence, since it reflects direct human preference) and is exempt from the
standard GRPO extraction pipeline.

\subsection{Loop 3: Reflection}

At the end of each session, Loop 3 produces a session summary:
\begin{itemize}[leftmargin=1.5em]
    \item Total tool calls, success rate, task completion rate.
    \item Top-3 recurring failure modes.
    \item Top-3 most-retrieved experiences (by session usage).
    \item Proposed new experiences (preliminary, before GRPO refinement).
\end{itemize}

This summary serves as the primary input to Agent-GRPO: the session
trajectories $\mathcal{S}$ fed to Algorithm~\ref{alg:agent-grpo} are
precisely the trajectories recorded during this session. Loop 3 fires
Agent-GRPO asynchronously after session end, so experience library updates
are available for the next session.

\subsection{Loop 4: Maintenance}

Loop 4 performs periodic library maintenance using statistics accumulated
across sessions:

\begin{itemize}[leftmargin=1.5em]
    \item \textbf{Staleness audit:} Experiences with $u = 0$ after
          $K_{\text{stale}} = 20$ sessions are flagged for deletion.
    \item \textbf{Confidence audit:} Experiences with $c < c_{\text{min}}$
          after at least $u \geq 5$ uses are deleted.
    \item \textbf{Redundancy audit:} Experience pairs with cosine similarity
          above $\theta_{\text{dup}}$ are merged via LLM consolidation.
    \item \textbf{Compression trigger:} When $|\mathcal{E}| > N_{\max}$,
          Loop 4 triggers the hybrid compression strategy
          (Section~\ref{sec:compression}).
\end{itemize}

The maintenance interval is configurable (default: every 10 sessions).

\subsection{Closed-Loop Diagram}

The integration of all four loops with Agent-GRPO forms a closed improvement
cycle:
\begin{equation}
    \underbrace{\text{Loop 0}}_{\text{Telemetry}}
    \;\xrightarrow{\tau}\;
    \underbrace{\text{Loop 3}}_{\text{Reflection}}
    \;\xrightarrow{\mathcal{S}}\;
    \underbrace{\text{Agent-GRPO}}_{\text{Alg.~\ref{alg:agent-grpo}}}
    \;\xrightarrow{\Omega^*}\;
    \underbrace{\mathcal{E}}_{\text{Library}}
    \;\xrightarrow{\text{top-}K}\;
    \underbrace{\text{Agent}}_{\text{next session}}
    \;\circlearrowleft
    \label{eq:closed-loop}
\end{equation}
with Loop 2 (user corrections) providing an out-of-band high-confidence
injection path and Loop 4 performing periodic quality control.

% ============================================================================
% 7. THEORETICAL ANALYSIS
% ============================================================================
